\chapter{Разработка концепции сервиса коллективных музыкальных рекомендаций для заведений и мероприятий}
\label{ch:service-concept}

\section{Описание работы сервиса}
\label{sec:service:concept}

В рамках настоящей работы предлагается концепция сервиса коллективных музыкальных рекомендаций для заведений и мероприятий на базе «Яндекс Музыки». Его назначение состоит в том, чтобы формировать и динамически обновлять музыкальный поток публичного воспроизведения с учетом предпочтений посетителей, фактически находящихся в пространстве заведения или мероприятия, при сохранении управляемости со стороны бизнеса. Иными словами, предметом рассматриваемой концепции является не индивидуальный выбор трека отдельным гостем, а управление общей музыкальной средой физического пространства, где необходимо одновременно учитывать вкусы группы, контекст площадки и ограничения владельца.

Необходимость такой концепции следует из рыночного разрыва, выявленного в первой главе. Существующие решения для музыки в бизнесе в основном помогают легально воспроизводить каталог, подбирать заранее подготовленные плейлисты, задавать жанровую рамку или управлять аудиомаркетингом. Эти функции важны, но они не решают задачу адаптации музыкального сопровождения к составу аудитории здесь и сейчас. В обычной модели заведение либо использует статический плейлист, либо полагается на музыкального редактора, персонал или DJ. Во всех этих случаях музыкальная программа, как правило, строится на предположении о типовой аудитории, а не на фактических предпочтениях людей, находящихся в помещении в конкретный момент.

Проектируемый сервис предлагает переход от статического или заранее подготовленного музыкального оформления к динамической коллективной рекомендации для физического пространства. Такая постановка отличается и от классических B2B-плееров, и от сервисов заказа музыки гостями. В первом случае музыкальный поток в основном определяется бизнесом заранее; во втором отдельный посетитель получает возможность запросить конкретный трек. Рассматриваемая продуктовая концепция занимает промежуточное положение: посетители не управляют эфиром напрямую, но их музыкальные предпочтения учитываются в агрегированном виде при формировании общего потока. Это позволяет сохранить сценарий заведения и одновременно сделать музыкальную среду более чувствительной к текущей аудитории.

В системе можно выделить три основные стороны. Первая сторона~--- посетитель, использующий приложение «Яндекс Музыки». Он выступает источником музыкальных предпочтений, накопленных в пользовательском профиле: истории прослушивания, реакций, персональных рекомендаций и других внутренних сигналов платформы. При этом участие посетителя должно быть добровольным: приложение может учитывать его музыкальный профиль только при наличии информированного согласия и необходимых разрешений. Вторая сторона~--- бизнес-пользователь: заведение, площадка или организатор события. Он задает контекст пространства, подключает устройство воспроизведения, определяет допустимые рамки музыкального потока и получает итоговый результат в виде списка треков-кандидатов или динамической очереди воспроизведения. Третья сторона~--- платформа «Яндекс Музыки», которая выступает вычислительным и рекомендательным контуром: сопоставляет сигналы присутствия, формирует текущую группу пользователей, учитывает настройки бизнес-профиля и генерирует музыкальный результат.

Базовая идея сервиса продолжает пользовательскую логику персональной «Моей Волны», но переносит ее с индивидуального слушателя на группу людей, совместно находящихся в одном физическом пространстве. Если у отдельных пользователей «Яндекс Музыки» уже есть история прослушивания, лайки, дизлайки и персональные рекомендации, то при корректном определении состава аудитории можно сформировать групповой профиль предпочтений и использовать его для генерации списка треков-кандидатов. Однако такая задача не сводится к механическому объединению индивидуальных вкусов. В группе могут присутствовать пользователи с разными музыкальными привычками, разной степенью вовлеченности и разной релевантностью для конкретного сценария заведения. Поэтому групповой музыкальный профиль должен рассматриваться как продуктовая абстракция, которая балансирует между множественностью предпочтений и необходимостью сформировать поток, приемлемый для пространства в целом.

В продуктовой логике это означает, что рекомендательная система не должна напрямую превращать все индивидуальные предпочтения посетителей в эфир. Сначала формируется множество треков-кандидатов, потенциально релевантных текущей группе, затем оно проходит через ограничения, заданные бизнесом. Такой подход особенно важен для заведений, где музыка является частью бренда и сценария обслуживания. Например, бар может допускать широкий диапазон популярных треков, но исключать слишком агрессивные жанры; ресторан может задавать более мягкую энергичность и ограничения по громкости или настроению; организатор вечеринки может использовать более динамичный режим, но сохранять стоп-листы нежелательных исполнителей. Детальные алгоритмы построения групповых рекомендаций рассматриваются в последующих технических главах, тогда как в данном подразделе важна сама продуктовая рамка: музыкальный поток формируется на пересечении предпочтений присутствующей аудитории и управляемого контекста площадки.

Для владельцев бизнеса, концертных площадок и организаторов мероприятий предполагается специальный профиль «Яндекс Музыка для бизнеса». Его роль шире, чем роль обычного плеера. Через этот профиль бизнес создает объект или мероприятие, привязывает к нему устройство воспроизведения, настраивает способ фиксации присутствия, выбирает режим рекомендаций и задает ограничения музыкальной среды. Эти ограничения могут включать жанровые рамки, уровень энергичности, настроение, временные режимы, стоп-листы треков или исполнителей, запрет на слишком частые повторения, ограничения по ненормативной лексике, допустимую степень популярности музыки и другие параметры, связанные с форматом площадки. Итоговая рекомендация строится не вместо этих ограничений, а внутри них: сервис не забирает у бизнеса контроль над атмосферой, а предлагает более гибкий способ наполнить заданные рамки треками, которые лучше соответствуют текущей аудитории.

Общая логика функционирования сервиса в физическом пространстве представлена на Рисунке~\ref{fig:service:general-scheme}. Схема показывает не алгоритмическую реализацию, а пространственный уровень концепции: в заведении есть устройство воспроизведения, связанное с бизнес-профилем, посетители со смартфонами и приложением «Яндекс Музыки», а также BLE-маяки, с помощью которых можно определить факт нахождения пользователя в определенной зоне.

\begin{figure}[htbp]
\centering
\includegraphics[width=0.98\textwidth]{service_general_scheme.png}
\caption{Общая схема функционирования сервиса коллективных музыкальных рекомендаций в пространстве заведения}
\label{fig:service:general-scheme}
\end{figure}

Как показано на Рисунке~\ref{fig:service:general-scheme}, проектируемый сервис связывает музыкальный поток с фактическим пространством заведения. Посетители находятся в баре, на танцполе, у входа или в других зонах и могут попадать в зону действия BLE-маяков. Устройство воспроизведения управляется через бизнес-профиль, а музыкальный поток формируется не по абстрактной аудитории, описанной заранее, а с учетом текущего состава посетителей, которые добровольно участвуют в сценарии. Такая схема особенно важна для баров, клубов и мероприятий, где аудитория меняется в течение вечера, а музыкальная среда должна адаптироваться к этой динамике.

Ключевым техническим предположением концепции является фиксация присутствия пользователя в заведении. Основной сценарий опирается на BLE-beacon: в пространстве размещается Bluetooth-маяк, передающий идентификатор ближайшим устройствам. Смартфон пользователя с установленным приложением «Яндекс Музыки» может зафиксировать нахождение рядом с таким маяком и, при наличии необходимых разрешений и информированного согласия, отправить на сервер сигнал о присутствии в конкретном объекте. В этой концепции идентификатор маяка не является источником персональных данных сам по себе; он используется как способ связать пользовательское событие присутствия с определенной площадкой или зоной внутри нее. Динамика входа и выхода посетителей, а также перемещение между зонами могут приводить к обновлению состава группы и, следовательно, к пересчету музыкальных рекомендаций.

Использование BLE не следует понимать как точную геолокацию пользователя внутри помещения. В продуктовой концепции речь идет о приближенном определении факта нахождения рядом с маяком или в зоне объекта, достаточном для формирования группового сценария. Такой подход позволяет автоматизировать участие пользователя и не требовать от него постоянных действий, но одновременно предъявляет повышенные требования к прозрачности коммуникации, разрешениям приложения и возможности отказаться от участия. Пользователь должен понимать, какие данные используются, зачем они нужны и почему владелец заведения не получает его индивидуальную музыкальную историю.

В качестве дополнительного или стартового сценария может использоваться QR-код. В этом случае посетитель сканирует код на площадке и явно присоединяется к музыкальному пространству. QR-механика проще для пилотного запуска, мероприятий или пространств, где BLE-инфраструктура еще не развернута: она не требует фоновой фиксации Bluetooth-сигнала и хорошо объясняется пользователю через видимый интерфейсный жест. Ее ограничение состоит в том, что пользователь должен совершить дополнительное действие, а состав группы может обновляться менее автоматически. Поэтому QR-код целесообразно рассматривать как альтернативный или резервный способ подключения, тогда как основной сценарий концепции остается связанным с BLE-маяками.

Информационный контур взаимодействия участников показан на Рисунке~\ref{fig:service:data-flow-scheme}. Если предыдущая схема описывала физическое пространство, то здесь отражена логика обмена данными между пользовательским приложением, серверной частью и бизнес-профилем сервиса.

\begin{figure}[htbp]
\centering
\includegraphics[width=0.98\textwidth]{service_data_flow_scheme.png}
\caption{Логика обмена данными между пользовательским приложением, серверной частью и бизнес-профилем сервиса}
\label{fig:service:data-flow-scheme}
\end{figure}

На Рисунке~\ref{fig:service:data-flow-scheme} серверный контур получает от приложения пользователя идентификатор пользователя и идентификатор BLE-маяка, сопоставляет маяк с бизнес-аккаунтом, формирует текущую группу пользователей и учитывает настройки, заданные со стороны бизнеса. После этого система формирует список треков-кандидатов и возвращает его в контур воспроизведения. Бизнес-профиль в этой логике выполняет двойную функцию: он является источником ограничений и одновременно получателем музыкального результата. Пользовательское приложение, в свою очередь, не передает заведению индивидуальную историю прослушивания, а участвует в формировании группового профиля через платформу.

Архитектурно сервис можно описать как последовательность операций. Сначала бизнес создает профиль объекта и задает ограничения. Затем пользователь, оказавшись в заведении, получает запрос на участие и дает информированное согласие. После этого платформа сопоставляет сигнал присутствия с конкретным объектом, формирует текущую группу пользователей, извлекает их музыкальные представления из внутреннего контура «Яндекс Музыки», строит групповой список кандидатов и применяет бизнес-ограничения. Результат передается в плеер бизнес-профиля и обновляется по мере изменения состава аудитории. В таком виде сервис соединяет пользовательский контур, B2B-контур и рекомендательный контур в единую продуктовую систему.

Такая концепция сознательно отделяет уровень рекомендаций от уровня промышленного лицензирования контента, но не игнорирует правовую рамку. Как показано в разделе~\ref{sec:market:legal}, для публичного воспроизведения музыки необходима отдельная правовая модель, включающая взаимодействие с правообладателями и организациями коллективного управления, а также B2B-урегулирование использования каталога. Поэтому проектируемый сервис следует рассматривать не просто как рекомендательный алгоритм, а как продуктовую систему, в которой рекомендательный слой, правовой слой, отчетность и пользовательское согласие должны быть согласованы. В данном подразделе правовая природа публичного воспроизведения не раскрывается повторно; важно лишь зафиксировать, что продуктовая концепция не может быть реализована изолированно от этих условий.

Таким образом, рассматриваемая концепция связывает несколько уровней: физическое пространство заведения, добровольное пользовательское участие, рекомендательный механизм и бизнес-управление музыкальной средой. Ее ключевая идея состоит в том, что музыкальный поток может формироваться не для условной усредненной аудитории, а для реально присутствующей группы людей, но при этом оставаться в рамках сценария, выбранного заведением или организатором. Далее в главе подробнее рассматриваются целевые аудитории и ценностное предложение, функции сервиса, пользовательские сценарии, а также позиционирование и наименование продуктовой концепции.

\section{Целевые аудитории сервиса и ценностное предложение}
\label{sec:service:audiences}

У сервиса есть четыре ключевые аудитории: владельцы и управляющие заведений, организаторы мероприятий, посетители и сама экосистема Яндекса. Для каждой аудитории ценность формулируется по-разному. Бизнесу важны управляемость музыкальной атмосферы, снижение операционной нагрузки и правовая определенность. Организатору мероприятия важен более точный «разогрев» аудитории и дополнительный сигнал о ее музыкальных предпочтениях. Посетителю важна возможность мягко влиять на музыкальный фон без необходимости специально заказывать треки или раскрывать свои данные заведению. Для Яндекса сервис может усилить ценность «Яндекс Музыки», создать новый B2B-сценарий и расширить понимание контекста музыкального потребления.

\begin{table}[htbp]
\caption{Ценностное предложение сервиса для ключевых аудиторий}
\label{tab:value-proposition}
\centering
\small
\singlespacing
\begin{tabularx}{\textwidth}{p{0.22\textwidth}X X}
\toprule
Аудитория & Проблема / потребность & Ценность сервиса \\
\midrule
Владелец заведения & Музыка зависит от вкуса персонала, плейлисты устаревают, правовая модель сложна & Управляемый музыкальный поток, ограничения по стилю, автоматизированная отчетность, более точное соответствие текущей аудитории \\
\addlinespace
Организатор мероприятия & Необходимо быстро вовлечь аудиторию и подобрать уместный стартовый музыкальный материал & Групповой профиль гостей, поддержка opening-сценария, возможность адаптировать музыку без ручного опроса аудитории \\
\addlinespace
Посетитель & Хочется влиять на атмосферу, но без навязчивого заказа треков и раскрытия личных данных площадке & Добровольное участие, влияние на общую «волну», сохранение персональных данных внутри контура Яндекса \\
\addlinespace
Яндекс & Требуется расширять экосистемные сценарии и ценность музыкального сервиса & Новый B2B-продукт, повышение вовлеченности, дополнительный контекст прослушивания, развитие технологий рекомендаций \\
\bottomrule
\end{tabularx}
\end{table}

Наиболее перспективными первыми сегментами выглядят бары, рестораны, кофейни, casual dining, event-friendly пространства, клубы и организаторы специальных мероприятий. В этих форматах музыка достаточно заметна для пользовательского опыта, но при этом не всегда требует полноценного живого DJ. Менее очевидными первыми сегментами являются медицинские учреждения, аптеки, часть retail-форматов и пространства с очень строгой брендовой музыкальной айдентикой. Для них стабильный контролируемый фон может быть важнее динамической персонализации.

Ценность сервиса для бизнеса не следует формулировать как прямое обещание роста выручки. Корректнее говорить о повышении управляемости клиентского опыта, снижении риска неуместного музыкального оформления, поддержке лояльности и более точном соответствии атмосферы аудитории. Экономический эффект может проявляться через длительность пребывания, повторные визиты, удовлетворенность и вовлеченность, но эти связи требуют отдельной эмпирической проверки.

\section{Функции сервиса для бизнеса и пользователей Яндекс Музыки}
\label{sec:service:functions}

Функциональность для бизнеса строится вокруг специального профиля «Яндекс Музыка для бизнеса». В этом профиле владелец объекта или организатор мероприятия создает карточку площадки, выбирает тип пространства, указывает режим использования музыки, подключает устройство воспроизведения и задает рамки рекомендаций. Базовые ограничения могут включать жанры, уровень энергичности, допустимый диапазон темпа, нежелательные исполнители, стоп-листы, запрет на повторения, ограничения по ненормативной лексике и сценарии по времени суток.

Отдельной функцией может стать AI-консультант для настройки музыкальных ограничений естественным языком. Вместо ручного выбора десятков параметров владелец может описать задачу: «нужна спокойная музыка для кофейни утром, без тяжелого рока и слишком популярных танцевальных треков» или «нужен мягкий разогрев для вечеринки в баре, без резких переходов и с постепенным повышением энергии». Система переводит такое описание в набор ограничений, который затем применяется к рекомендательному потоку. В рамках настоящей ВКР этот элемент рассматривается как продуктовая возможность, а не как реализуемая техническая часть эксперимента.

Для бизнеса также важна отчетность. Сервис должен фиксировать фактические воспроизведения, хранить идентификаторы треков, исполнителей, авторов и изготовителей фонограмм в объеме, необходимом для подготовки отчетов РАО и ВОИС. Автоматическая генерация отчетов снижает административную нагрузку и делает продукт более ценным в сравнении с простым рекомендательным плейлистом. При этом конкретный формат отчетных файлов должен соответствовать актуальным требованиям организаций по коллективному управлению правами~\cite{rao_report_forms,vois_report_form}.

Функциональность для посетителя должна быть максимально ненавязчивой. В обычном сценарии пользователь не открывает отдельное приложение и не выбирает треки вручную. Если приложение фиксирует нахождение пользователя в пространстве, где используется сервис, пользователь получает понятное уведомление: он может подключиться к музыкальному пространству и разрешить использовать свои музыкальные предпочтения для формирования общей волны. При отказе пользователь не участвует в рекомендациях, а его пребывание не должно влиять на музыкальный поток.

Принципиальное требование к пользовательскому сценарию~--- сохранение приватности. Владелец заведения не видит список присутствующих пользователей, их историю прослушивания, лайки, дизлайки или индивидуальные музыкальные профили. Он получает только результат в виде плейлиста, очереди воспроизведения и агрегированных операционных данных. Обработка персональных данных остается внутри контура Яндекса, а пользовательское согласие должно быть явно получено и отзываемо.

\section{Пользовательские сценарии взаимодействия с сервисом}
\label{sec:service:journeys}

Сценарий владельца заведения начинается с подключения бизнес-профиля. Пользователь со стороны бизнеса выбирает тип объекта, указывает адрес, подключает устройство воспроизведения, выбирает способ фиксации присутствия посетителей и задает начальные музыкальные рамки. На старте сервис может предложить шаблон для ресторана, бара, кофейни или мероприятия, а затем уточнить параметры: уровень энергии, допустимые жанры, нежелательные направления, язык треков, режим обновления и необходимость отчетности.

В операционном режиме владелец или администратор не должен постоянно управлять музыкой вручную. Сервис формирует очередь, обновляет ее при изменении состава аудитории и показывает базовые индикаторы: активен ли маяк, сколько пользователей участвует в формировании рекомендаций в агрегированном виде, не нарушены ли ограничения, сформирован ли отчет. При необходимости администратор может временно зафиксировать плейлист, усилить или снизить энергичность, заблокировать отдельный трек или переключиться на заранее подготовленный fallback-плейлист.

Сценарий организатора мероприятия похож, но имеет более выраженную временную структуру. До события организатор создает карточку мероприятия, задает музыкальный профиль, подключает QR-код или BLE-зону и может заранее пригласить гостей подключиться к музыкальному пространству. В начале события сервис формирует мягкий opening-поток, затем постепенно адаптируется к фактическому составу гостей. В кульминационной части мероприятия система может работать как источник рекомендаций для DJ или как автоматический слой для фоновых зон, не претендуя на замену живого музыкального управления.

Сценарий посетителя должен быть простым. Пользователь приходит в заведение или на мероприятие, получает уведомление или видит QR-код, знакомится с коротким объяснением и принимает или отклоняет участие. Если он соглашается, его музыкальные предпочтения учитываются в агрегированном виде. Пользователь не выбирает конкретный трек и не получает индивидуального контроля над эфиром; его вклад становится частью коллективной волны. Такая модель снижает риск конфликтов, характерный для сервисов одиночного заказа музыки, где один активный или платежеспособный гость может слишком сильно влиять на атмосферу.

Отдельный сценарий связан с использованием технологии как помощника DJ. DJ может загрузить или сопоставить подготовленный плейлист и получить оценку его потенциального соответствия аудитории, находящейся на площадке. В более развитой версии система могла бы анализировать музыкальные файлы по аудиоэмбеддингам и сопоставлять их с групповым профилем. В настоящей работе этот сценарий рассматривается как перспективное применение технологии, но не реализуется экспериментально.

Еще одно возможное применение связано с анализом вовлеченности аудитории по зонам пространства. При наличии нескольких BLE-меток теоретически можно различать посетителей, находящихся у бара, в зале или ближе к танцполу, и адаптировать музыкальный поток с учетом пространственной структуры аудитории. Такой сценарий требует особенно аккуратной работы с приватностью и не входит в техническую реализацию ВКР, но может быть полезен для дальнейшего развития продукта.

\section{Позиционирование, наименование и коммуникационная концепция сервиса}
\label{sec:service:positioning}

После описания продуктовой логики, целевых аудиторий, функций и пользовательских сценариев необходимо определить коммуникационную рамку сервиса. В рамках настоящей работы проектируемое направление целесообразно обозначать как сервис «Яндекс Музыка для бизнеса», а «Резонанс» рассматривать как название функции коллективной адаптации музыкального потока. Эти наименования не являются окончательным коммерческим решением, но используются в работе как концептуальная рамка, позволяющая связать техническую, пользовательскую и маркетинговую логику продукта.

Выбор названия должен учитывать двойственную природу продукта. С одной стороны, сервис относится к музыкальной среде и должен ассоциироваться со звучанием, совместным восприятием и атмосферой пространства. С другой стороны, он не является обычным музыкальным плеером или сервисом заказа треков: его ключевая идея состоит в согласовании индивидуальных музыкальных предпочтений нескольких посетителей в общий поток, приемлемый для заведения или мероприятия. Поэтому название должно передавать не столько факт воспроизведения музыки, сколько эффект совпадения, настройки и совместного звучания.

С этой точки зрения слово «Резонанс» представляется содержательно обоснованным для пользовательской функции сервиса. В физическом смысле резонанс связан с усилением колебаний при совпадении частот. В коммуникационном смысле эта метафора может быть перенесена на согласование вкусов разных людей: музыкальные предпочтения посетителей не просто суммируются, а приводятся к общему звучанию, которое должно соответствовать текущей аудитории и контексту площадки. Таким образом, название отражает центральную идею функции: создание музыкальной среды, в которой индивидуальные профили пользователей начинают работать как часть коллективного сценария.

Для связи с существующей экосистемой целесообразно использовать связку «Яндекс Музыка для бизнеса» как название сервисного направления и «Резонанс» как название функции коллективной адаптации музыкального потока. Она показывает, что продукт не создается как независимая внешняя платформа, а является возможным B2B/B2C-слоем внутри «Яндекс Музыки». Такая связка снижает необходимость объяснять пользователю источник музыкальных предпочтений и подчеркивает, что сервис опирается на уже существующую платформу, пользовательские профили и рекомендательные технологии. При этом в дальнейшей коммуникации возможны разные уровни использования названия: «Яндекс Музыка для бизнеса» для официального B2B-позиционирования и «Резонанс» для интерфейсных и пользовательских сценариев.

Пользовательская коммуникация может строиться вокруг формул «Войти в резонанс» или «Резонировать». Эти выражения удобны тем, что описывают участие пользователя не как техническую передачу данных, а как добровольное присоединение к музыкальному пространству. Однако такая метафора не должна заменять юридически значимое согласие. В интерфейсе необходимо отдельно и ясно объяснять, какие данные используются, с какой целью они обрабатываются и почему владелец заведения не получает доступ к индивидуальному музыкальному профилю посетителя. Иначе эмоционально удачная формула может создать риск недопонимания механики сервиса.

Для бизнеса позиционирование должно быть более прагматичным. Сервис «Яндекс Музыка для бизнеса» и функция «Резонанс» должны описываться не как инструмент гарантированного роста выручки, а как решение для управляемой музыкальной среды, которое помогает учитывать текущую аудиторию, сохранять рамки бренда и снижать операционную сложность работы с музыкой. Возможная формула ценности для владельца заведения может быть сформулирована следующим образом: «музыка, которая остается вашей по стилю, но лучше совпадает с гостями, находящимися здесь сейчас». Такая формулировка подчеркивает, что сервис не отменяет музыкальную политику площадки, а работает внутри заданных бизнесом ограничений.

Для организаторов мероприятий функция «Резонанс» может позиционироваться как инструмент мягкого вовлечения аудитории и поддержки музыкальной драматургии события. В отличие от голосований и сервисов заказа треков, она не создает прямой конкуренции посетителей за эфир и не передает управление музыкальной программой одному наиболее активному гостю. Сервис работает с агрегированным профилем группы, поэтому может быть полезен в ситуациях, где музыка должна поддерживать общий сценарий вечера и при этом учитывать фактический состав аудитории.

Отдельным элементом коммуникационной концепции должна стать приватность. Название «Резонанс» и связанные с ним пользовательские формулы работают только при условии доверия к механике сервиса. Пользователю необходимо явно сообщить, что его история прослушивания не передается заведению, не показывается администратору и не используется для идентификации перед третьими лицами. Бизнес получает музыкальный результат, очередь воспроизведения и агрегированные показатели, но не персональные вкусовые профили гостей. Такая коммуникация помогает отличить сервис от технологий наблюдения за посетителями и делает участие пользователя более осознанным.

Таким образом, предлагаемое в настоящей работе наименование функции «Резонанс» в рамках сервиса «Яндекс Музыка для бизнеса» выполняет не только маркирующую, но и объясняющую функцию. Оно связывает музыкальную природу продукта, идею коллективного согласования предпочтений и сценарий физического пространства. В дальнейшем это название может быть уточнено в ходе пользовательских исследований и брендовой экспертизы, однако на уровне концепции оно позволяет последовательно описывать сервис как продукт, в котором музыкальная среда заведения настраивается на фактическую аудиторию, не разрушая контроль со стороны бизнеса и требования приватности.

\section{Модель монетизации, дистрибуции и экономические эффекты внедрения сервиса}
\label{sec:service:monetization}

\subsection{Подход к монетизации и тарифная модель}
\label{subsec:service:monetization:tariffs}

Обновленная оценка рынка в разделе~\ref{sec:market:potential} показывает, что проектируемый сервис должен рассматриваться не как массовое решение для всех заведений, где в принципе звучит музыка, а как продукт для ограниченного сегмента музыкально чувствительных форматов. В рамках этой оценки TAM составляет около 14--23 тыс. точек, SAM~--- около 3--5 тыс. точек, а реалистично достижимый SOM на горизонте раннего коммерческого масштабирования~--- 450--550 заведений при базовом ориентире около 500 точек. Следовательно, модель монетизации должна быть достаточно простой для B2B-продаж и одновременно соответствовать повышенной ценности продукта: сервис предлагает не только легальный музыкальный плеер или набор плейлистов, а динамическую адаптацию музыкальной среды к фактическому составу аудитории.

Основной моделью монетизации целесообразно считать B2B-подписку за подключенную точку или объект. Такая модель лучше всего согласуется с постоянным рынком баров, клубов, restaurant-bar и lounge-форматов, выделенным в разделе~\ref{sec:market:potential}: заведение регулярно использует музыкальное сопровождение, получает повторяющуюся ценность и может планировать расходы как часть операционного бюджета. В качестве предварительной ценовой гипотезы в настоящей работе используется вилка 6\,000--9\,000~руб. в месяц за одну точку. Эта цена не является финально утвержденным коммерческим тарифом; она нужна для проверки жизнеспособности MVP и сопоставления потенциальной выручки с достижимым SOM.

Выбор такой вилки объясняется несколькими соображениями. Во-первых, цена должна быть заметно ниже совокупных платежей за публичное воспроизведение музыки для типового релевантного заведения, рассмотренных в разделе~\ref{sec:market:potential}, чтобы сервис воспринимался как дополнительный технологический слой, а не как замена лицензионным обязательствам перед РАО и ВОИС. Во-вторых, цена должна быть выше уровня символического SaaS-платежа, поскольку продукт требует серверного расчета рекомендаций, интеграции с пользовательскими профилями, настройки бизнес-ограничений и поддержки B2B-клиента. В-третьих, сервис претендует на дифференциацию: в отличие от стандартных решений фоновой музыки, он соединяет широкую пользовательскую базу «Яндекс Музыки», данные о предпочтениях слушателей, рекомендательные технологии и офлайн-контекст заведения.

Предварительная тарифная логика может быть построена по трехуровневой модели, представленной в Таблице~\ref{tab:service:tariffs}. Она не фиксирует окончательные цены и состав пакетов, но показывает, как ценность сервиса может быть разложена по уровням сложности клиента.

\begin{table}[htbp]
\centering
\small
\caption{Предварительная тарифная логика сервиса}
\label{tab:service:tariffs}
\begin{tabularx}{\textwidth}{p{0.20\textwidth}p{0.23\textwidth}X}
\toprule
Тарифный уровень & Целевой клиент & Состав ценности \\
\midrule
Базовый & Одна точка с одним основным музыкальным контуром & Коллективная рекомендация на основе присутствующих пользователей, базовые фильтры по стилю и энергичности, ограниченная отчетность, стандартная поддержка. \\
\addlinespace
Расширенный & Заведение с несколькими режимами, зонами или регулярными событиями & Несколько зон или сценариев внутри объекта, более глубокая настройка параметров музыкальной среды, режимы по времени суток и формату мероприятий, расширенная аналитика, более высокая частота обновления музыкального потока, приоритетная поддержка. \\
\addlinespace
Сетевой / корпоративный & Сеть заведений, управляющая компания или крупный B2B-клиент & Индивидуальный расчет, несколько объектов, централизованное управление музыкальными политиками, единая отчетность, возможная интеграция с существующими B2B-продуктами экосистемы. \\
\bottomrule
\end{tabularx}
\end{table}

В этой логике базовый тариф отвечает за проверку основной продуктовой гипотезы: готов ли бизнес платить за адаптивную музыкальную среду, зависящую от текущей аудитории. Расширенный тариф позволяет монетизировать более сложные сценарии, где заведение хочет управлять разными зонами, временными режимами и форматами вечеров. Сетевой тариф нужен для клиентов, у которых один договор может охватывать несколько физических точек, что особенно важно для достижения SOM: подключение сетевых клиентов сокращает стоимость продаж на одну точку и делает раннее масштабирование более реалистичным.

Событийная модель монетизации также может быть предусмотрена, но в данном подразделе она не является основой финансового расчета. Для мероприятий возможны разовый тариф за событие, дневной доступ или расширение подписки для баров и клубов, проводящих регулярные вечеринки. Однако события отличаются единицей учета, бюджетодержателем и циклом продажи, поэтому их целесообразно рассматривать как дополнительный канал применения сервиса. Расчет окупаемости MVP далее строится именно на подписке постоянных B2B-точек, поскольку эта модель напрямую связана с оценкой SOM из раздела~\ref{sec:market:potential}.

\subsection{Каналы дистрибуции и целевые сценарии продаж}
\label{subsec:service:monetization:distribution}

Дистрибуция сервиса должна соответствовать тому, что продукт находится на пересечении потребительской платформы «Яндекс Музыки» и B2B-рынка заведений. Наиболее сильным каналом может стать экосистемная дистрибуция Яндекса. Она предполагает продвижение через уже существующие направления работы с бизнесом, использование партнерской базы, кросс-продажи с близкими продуктами и постепенное встраивание сервиса в знакомую для клиентов инфраструктуру. Для заведения ценность такого канала состоит в снижении неопределенности: продукт воспринимается не как отдельный неизвестный стартап, а как расширение платформы, с которой уже связаны пользовательские музыкальные профили, устройства воспроизведения и рекомендательные технологии.

Возможная экосистемная дистрибуция включает несколько уровней. На уровне «Яндекс Музыки» сервис может предлагаться как B2B-режим для заведений, которым важна легальная и управляемая музыкальная среда. На уровне устройств и голосовых интерфейсов возможна связь с уже привычными сценариями воспроизведения через «Алису» и устройства Яндекса. На уровне продуктов для бизнеса сервис может быть предложен как дополнительный модуль для заведений, которые уже используют решения экосистемы для ресторанов, малого бизнеса, рекламы, навигации или клиентской коммуникации. Такая логика не означает автоматического доступа ко всем этим каналам, но задает направление коммерческого развития: продукт выигрывает от того, что его не нужно выводить на рынок в полной изоляции от существующей экосистемы.

Второй канал~--- прямые B2B-продажи в наиболее релевантные сегменты. В фокусе должны находиться бары, клубы, restaurant-bar, lounge-форматы, площадки с танцполом и музыкально ориентированные заведения, выделенные в первой главе. Для таких клиентов музыка не является случайным фоном; она влияет на восприятие вечера, длительность пребывания, вовлеченность гостей и соответствие атмосферы позиционированию заведения. Поэтому аргументация продажи должна строиться не вокруг абстрактной «инновационности» сервиса, а вокруг управляемого улучшения аудиосреды: бизнес сохраняет контроль над стилем, но получает возможность учитывать фактический состав аудитории.

Третий канал связан с пилотными внедрениями. Для нового B2B-продукта недостаточно показать расчетную экономику; необходимо проверить, как сервис работает в реальных заведениях, насколько понятно бизнесу ценностное предложение, готовы ли пользователи давать согласие на участие и не создает ли технология операционных сложностей для персонала. Поэтому ранний запуск должен включать ограниченное число площадок, сбор обратной связи, проверку готовности платить, оценку стабильности обновления рекомендаций и формирование кейсов для дальнейших продаж. В рамках пилота особенно важно не обещать доказанный рост выручки, а фиксировать более проверяемые показатели: стабильность работы, долю пользователей, согласившихся участвовать, частоту обновления группы, удовлетворенность персонала и субъективную оценку музыкального соответствия.

Четвертый канал~--- партнерский. Потенциальными партнерами могут выступать интеграторы оборудования, поставщики решений для HoReCa, операторы цифровой инфраструктуры в заведениях, B2B-поставщики музыки, event-платформы и компании, уже имеющие физическую инфраструктуру в общественных пространствах. Партнерская модель особенно важна для снижения затрат на вход: если часть оборудования, каналов продаж или клиентских отношений уже существует, сервису не нужно заново строить весь контур присутствия в заведениях.

Отдельно следует рассмотреть временную технологическую стратегию пилота. В качестве проектного допущения на первом этапе можно использовать уже размещенную инфраструктуру в заведениях как потенциальный способ снизить капитальные расходы на разворачивание собственной сети BLE-маяков. В частности, Яндекс в 2024 году сообщил о покупке сервиса аренды зарядных устройств «Бери заряд!» и развитии арендной инфраструктуры как отдельного направления Яндекс Go~\cite{yandex_beri_zaryad_2024}. Это не доказывает техническую готовность таких станций выполнять BLE-сценарий и не означает, что они могут быть использованы без дополнительной разработки и договорной проработки. Однако сам факт наличия размещаемой в заведениях аппаратной инфраструктуры позволяет рассматривать подобные объекты как один из возможных пилотных контуров: при наличии технической совместимости и правовых оснований они могли бы уменьшить потребность в самостоятельной закупке и монтаже маяков на ранней стадии.

\subsection{Минимальная конфигурация MVP и оценка затрат на запуск}
\label{subsec:service:monetization:mvp}

Экономическая оценка сервиса должна исходить из того, что MVP создается не как самостоятельная музыкальная платформа с нуля, а как надстройка над уже существующими активами Яндекса. В минимальной конфигурации предполагается использование инфраструктуры «Яндекс Музыки», пользовательских профилей, серверного контура, авторизации, рекомендательных технологий и близких по смыслу B2B-инструментов. Такое допущение принципиально снижает стартовую смету: в расчет не включается создание каталога, полноценного стримингового сервиса, собственной системы аккаунтов, отдельной платежной инфраструктуры и полного контура воспроизведения с нуля.

Минимальный MVP должен проверить не всю будущую продуктовую концепцию, а ее экономически значимое ядро. К нему относятся: подключение ограниченного числа B2B-точек, создание бизнес-профиля или близкого административного интерфейса, получение согласия пользователя на участие в сценарии, фиксация присутствия через BLE или QR, формирование текущей группы пользователей, расчет групповой рекомендации, применение базовых бизнес-фильтров и выдача результата в контур воспроизведения. Расширенная аналитика, сложное много-зонное управление, глубокая автоматизация отчетности и сетевые кабинеты могут быть перенесены за пределы первого MVP, если они не требуются для проверки готовности платить.

Сценарный горизонт разработки MVP можно оценить в 4--6 месяцев. Команда должна быть компактной, но межфункциональной: продуктовый менеджер или product owner, UX/UI-дизайнер частичной занятости, backend-инженеры, мобильный разработчик для изменений в пользовательском приложении, ML/recommendation engineer для адаптации групповой модели к продуктовой логике, аналитик данных, QA-инженер, DevOps или инфраструктурная поддержка частичной занятости, юрист или privacy specialist для согласий и обработки данных, а также B2B sales/account manager на этапе пилота. Такая команда не равна постоянному штату отдельного бизнеса; в крупной экосистеме часть ролей может подключаться долями занятости.

Предварительная оценка минимальных затрат представлена в Таблице~\ref{tab:service:mvp-cost}. Она является расчетом автора на основе принятых допущений о составе команды, длительности работ и использовании существующей инфраструктуры. Цель оценки~--- определить порядок затрат для проверки окупаемости, а не построить детальный бюджет промышленного запуска.

\begin{table}[htbp]
\centering
\small
\caption{Предварительная оценка затрат на разработку и запуск MVP}
\label{tab:service:mvp-cost}
\begin{tabularx}{\textwidth}{p{0.28\textwidth}p{0.20\textwidth}X}
\toprule
Блок затрат & Оценка & Комментарий \\
\midrule
Продукт, дизайн и аналитика & 3--5 млн~руб. & Product owner, UX/UI частичной занятости, аналитик данных, подготовка пилотных сценариев и метрик. \\
\addlinespace
Backend, мобильная интеграция и рекомендательный контур & 11--15 млн~руб. & Разработка серверной логики, интеграции с пользовательским приложением, адаптация модели групповых рекомендаций, базовые фильтры и бизнес-профиль. \\
\addlinespace
QA, инфраструктура и эксплуатационная подготовка & 3--4 млн~руб. & Тестирование, мониторинг, DevOps-поддержка, подготовка ограниченного пилотного контура. \\
\addlinespace
Правовая и privacy-проработка & 1--2 млн~руб. & Проверка согласий, сценариев обработки данных, B2B-договорной модели и ограничений использования инфраструктуры. \\
\addlinespace
Пилотные внедрения и резерв & 6--8 млн~руб. & Подключение первых площадок, настройка оборудования или партнерской инфраструктуры, поддержка B2B-клиентов, резерв на интеграционные риски. \\
\midrule
\textbf{Итого} & \textbf{24--32 млн~руб.} & Базовый ориентир для дальнейшего расчета~--- 28 млн~руб. \\
\bottomrule
\end{tabularx}
\end{table}

Эта смета является сдержанной именно потому, что MVP опирается на уже существующие активы. Если бы продукт создавался вне крупной платформы, потребовались бы существенно более высокие расходы на каталог, лицензирование, инфраструктуру воспроизведения, систему аккаунтов, рекомендательный стек и поддержку. В рассматриваемой модели основные затраты связаны не с созданием музыкального стриминга, а с интеграцией группового рекомендательного слоя в продуктовый и B2B-контур.

Отдельное значение имеет аппаратный слой фиксации присутствия. В базовую MVP-смету не закладывается массовая закупка и монтаж собственной сети BLE-маяков, поскольку на раннем этапе рациональнее использовать QR-сценарий, ограниченное число тестовых маяков или уже размещенную партнерскую инфраструктуру при наличии технической и договорной возможности. Такой подход не отменяет будущих расходов на оборудование при масштабировании, но позволяет проверить продуктовую гипотезу до того, как проект понесет значительные капитальные затраты.

\subsection{Предварительная unit-экономика и расчет точки безубыточности}
\label{subsec:service:monetization:unit}

Для проверки финансовой реалистичности MVP используется упрощенная модель, в которой основными затратами считаются фиксированные расходы на разработку и пилотный запуск. Переменные расходы на обслуживание одной точки на стадии MVP предполагаются сравнительно низкими благодаря использованию существующей инфраструктуры Яндекса; поэтому они не выделяются в отдельную расчетную строку. Такое упрощение допустимо для первичной оценки, но при переходе к промышленному масштабированию модель должна быть расширена: в нее необходимо включить поддержку клиентов, B2B-продажи, сопровождение оборудования, маркетинг, инфраструктурные расходы и юридическое сопровождение.

Годовая выручка от подписочной модели рассчитывается по формуле:

\begin{equation*}
    R_{\mathrm{year}} = N \cdot P_{\mathrm{sub}} \cdot 12,
\end{equation*}

где $R_{\mathrm{year}}$~--- годовая выручка, $N$~--- число подключенных точек, $P_{\mathrm{sub}}$~--- месячная цена подписки за одну точку. В качестве сценариев используются три цены: 6\,000~руб., 7\,500~руб. и 9\,000~руб. в месяц. Для числа подключенных точек берется диапазон SOM из раздела~\ref{sec:market:potential}: консервативный сценарий~--- 450 точек, базовый~--- 500 точек, оптимистичный~--- 550 точек.

\begin{table}[htbp]
\centering
\small
\caption{Годовая выручка при достижении SOM, млн~руб.}
\label{tab:service:som-revenue}
\begin{tabularx}{\textwidth}{p{0.26\textwidth}XXX}
\toprule
Цена подписки за точку & SOM 450 точек & SOM 500 точек & SOM 550 точек \\
\midrule
6\,000~руб./мес. & 32{,}4 & 36{,}0 & 39{,}6 \\
7\,500~руб./мес. & 40{,}5 & 45{,}0 & 49{,}5 \\
9\,000~руб./мес. & 48{,}6 & 54{,}0 & 59{,}4 \\
\bottomrule
\end{tabularx}
\end{table}

Как видно из Таблицы~\ref{tab:service:som-revenue}, достижение даже нижней границы SOM формирует годовую выручку от 32{,}4 до 48{,}6 млн~руб. в зависимости от цены подписки. При базовом SOM в 500 точек диапазон составляет 36--54 млн~руб. в год. Эти значения не следует трактовать как прогноз гарантированной выручки: они показывают, какой порядок дохода может возникнуть, если продукт подтвердит ценность и достигнет раннего коммерческого масштаба, обозначенного в первой главе.

Точка безубыточности по числу подключенных точек рассчитывается следующим образом:

\begin{equation*}
    N_{\mathrm{BE}} = \frac{C_{\mathrm{MVP}}}{P_{\mathrm{sub}} \cdot 12},
\end{equation*}

где $N_{\mathrm{BE}}$~--- количество точек, необходимое для покрытия затрат на MVP, $C_{\mathrm{MVP}}$~--- стоимость разработки и пилотного запуска, $P_{\mathrm{sub}}$~--- месячная цена подписки. Для оценки используется диапазон стоимости MVP 24--32 млн~руб. и базовый ориентир 28 млн~руб.

\begin{table}[htbp]
\centering
\small
\caption{Количество точек, необходимое для покрытия затрат на MVP}
\label{tab:service:break-even}
\begin{tabularx}{\textwidth}{p{0.22\textwidth}XXXX}
\toprule
Цена подписки & MVP 24 млн~руб. & MVP 28 млн~руб. & MVP 32 млн~руб. & Доля базового SOM при MVP 28 млн~руб. \\
\midrule
6\,000~руб./мес. & 334 & 389 & 445 & 78\% \\
7\,500~руб./мес. & 267 & 312 & 356 & 62\% \\
9\,000~руб./мес. & 223 & 260 & 297 & 52\% \\
\bottomrule
\end{tabularx}
\end{table}

Таблица~\ref{tab:service:break-even} показывает, что при базовой MVP-смете 28 млн~руб. сервису необходимо подключить около 389 точек при цене 6\,000~руб. в месяц, около 312 точек при цене 7\,500~руб. и около 260 точек при цене 9\,000~руб. В пересчете на базовый SOM в 500 точек это составляет соответственно 78\%, 62\% и 52\% достижимого раннего коммерческого сегмента. Даже в более жестком сценарии стоимости MVP 32 млн~руб. и минимальной цене 6\,000~руб. точка безубыточности составляет около 445 точек, то есть остается близкой к нижней границе SOM. При средней цене 7\,500~руб. и базовой смете 28 млн~руб. выход на безубыточность достигается до полного освоения консервативного SOM в 450 точек.

Полученный результат важен для оценки реалистичности запуска. Проект не требует охвата всего SAM или значительной доли TAM для покрытия расходов на MVP: при сдержанной смете и использовании инфраструктуры Яндекса достаточно подключить ограниченную часть рынка, уже выделенного в разделе~\ref{sec:market:potential}. При этом расчет остается предварительным. Он не учитывает возможные скидки для первых клиентов, длительность цикла продажи, отток, расходы на сопровождение сетевых клиентов, налоговые и бухгалтерские аспекты, а также дополнительные затраты на оборудование при расширении BLE-инфраструктуры. Поэтому вывод следует формулировать осторожно: модель показывает потенциальную экономическую проверяемость MVP, но не доказывает автоматическую коммерческую успешность сервиса.

\subsection{Экосистемные эффекты для Яндекса и потенциал использования данных}
\label{subsec:service:monetization:ecosystem}

Прямая подписочная выручка не исчерпывает ценность проектируемого сервиса. Для Яндекса продукт может выступать как способ усиления «Яндекс Музыки» за счет нового сценария использования: музыкальные предпочтения пользователя начинают влиять не только на его индивидуальное прослушивание, но и на атмосферу реального физического пространства. Это отличает сервис и от обычного B2C-стриминга, где рекомендации остаются внутри личного опыта, и от B2B-сервисов фоновой музыки, которые работают преимущественно с заранее подготовленными плейлистами и настройками бизнеса.

Такой сценарий может стать уникальным торговым предложением «Яндекс Музыки» на B2B-рынке. Конкурентное отличие возникает не из отдельной функции воспроизведения, а из сочетания нескольких активов: широкой базы пользователей, их музыкальных профилей, рекомендательных технологий, узнаваемого бренда и возможности связать онлайн-профиль слушателя с офлайн-контекстом заведения. Для бизнеса это может быть аргументом в пользу выбора решения на базе «Яндекс Музыки», поскольку стандартный музыкальный поставщик не обладает сопоставимой связью с пользовательскими предпочтениями посетителей.

Для пользователей сервис также создает дополнительный стимул к вовлеченности. Если человек понимает, что его музыкальный профиль может мягко влиять на атмосферу заведения или мероприятия, у него появляется причина сохранять приложение установленным, быть авторизованным, поддерживать актуальность профиля через лайки и прослушивания и добровольно участвовать в сценарии. Этот эффект не следует преувеличивать: далеко не каждый посетитель захочет подключаться к такому механизму. Однако для части аудитории ценность участия может быть понятнее, чем в абстрактной системе мониторинга, поскольку пользователь получает осязаемый результат в виде более релевантной музыкальной среды.

Еще один эффект связан с развитием B2B-направления экосистемы. Сервис соединяет потребительскую музыкальную платформу и бизнес-аудиторию заведений, создавая новый тип продукта на границе B2C и B2B. Это может поддерживать кросс-продажи с другими решениями Яндекса для бизнеса, усиливать ценность экосистемного пакета и расширять присутствие компании в офлайн-среде. При этом продуктовая логика не должна сводиться к навязыванию еще одного сервиса: устойчивость модели зависит от того, увидит ли бизнес самостоятельную ценность в адаптивной музыкальной среде, отчетности и управляемости музыкального сценария.

Отдельного внимания заслуживает потенциал данных, возникающих в таком продукте. Современный музыкальный стриминг в основном располагает информацией об индивидуальном цифровом поведении пользователя: прослушиваниях, лайках, пропусках, поисковых действиях и реакциях на рекомендации. Проектируемый сервис может добавить контекст группового и пространственного потребления: какие музыкальные предпочтения оказываются совместимыми в конкретных типах заведений, как меняется состав аудитории в течение вечера, какие настройки бизнеса чаще соответствуют фактическому отклику посетителей, какие временные слоты требуют более энергичного или более спокойного музыкального потока. В дальнейшем такие агрегированные данные могут использоваться для улучшения рекомендательных моделей и для создания B2B-аналитики.

Коммерческое использование данных в этой модели возможно только в правомерном, агрегированном и обезличенном виде. Нельзя рассматривать сервис как механизм передачи индивидуальных пользовательских профилей владельцам заведений или третьим лицам. Допустимыми направлениями могут быть улучшение внутренних моделей рекомендаций, аналитические отчеты для самих B2B-клиентов в рамках продукта, выявление паттернов музыкальной атмосферы по типам заведений, времени суток и формату событий, а также развитие платных аналитических модулей при соблюдении требований к защите данных и информированному согласию. Например, бизнесу могут быть доступны не личные вкусы конкретных гостей, а агрегированные показатели: как менялась музыкальная релевантность в течение вечера, какие режимы чаще выбирались для разных временных слотов, насколько активно пользователи соглашались участвовать в сценарии.

Таким образом, модель монетизации через подписку 6\,000--9\,000~руб. в месяц выглядит проверяемой на рынке, выделенном в разделе~\ref{sec:market:potential}. При оценке MVP в 24--32 млн~руб. и базовом ориентире 28 млн~руб. расчетная точка безубыточности достигается в пределах SOM 450--550 точек, особенно при среднем и расширенном ценовых сценариях. Помимо прямой выручки, сервис может создавать стратегические эффекты для Яндекса: дифференцировать «Яндекс Музыку», усилить B2B-направление, дать новый контекст для развития рекомендаций и открыть возможность для агрегированной аналитики музыкальной среды. Все эти эффекты должны рассматриваться как гипотезы дальнейшей проверки, а не как уже доказанный коммерческий результат.

\section{Выводы}
\label{sec:service:conclusion}

В данной главе разработана продуктовая концепция сервиса «Яндекс Музыка для бизнеса» и функции «Резонанс» как B2B/B2C-слоя внутри экосистемы «Яндекс Музыки». Сервис предполагает специальный бизнес-профиль для заведений и мероприятий, механизм определения присутствия посетителей через BLE или QR, добровольное участие пользователей, обработку данных внутри контура Яндекса и формирование динамического группового музыкального потока.

Ключевая продуктовая идея состоит в том, что заведение сохраняет контроль над музыкальной атмосферой, а посетители получают мягкое влияние на нее без индивидуального заказа треков и без передачи персональных данных площадке. Такое решение отличается от классических B2B-сервисов музыки для бизнеса тем, что учитывает фактический состав аудитории, и отличается от сервисов заказа музыки тем, что не дает одному гостю непропорционального влияния на эфир.

Функциональность сервиса должна включать настройку рамок воспроизведения, стоп-листы, управление энергичностью, возможный AI-консультант для задания ограничений естественным языком, автоматическую генерацию отчетности для РАО и ВОИС, а также сценарии для мероприятий и поддержки DJ. Монетизация может строиться отдельно для постоянных точек и событий, поскольку эти рынки различаются по циклу продажи, LTV и логике использования.

Сформулированная концепция задает бизнес- и продуктовый контекст для последующих технических глав. Далее работа переходит к теоретическим и технологическим основам музыкальных рекомендательных систем, подходам к групповым рекомендациям и экспериментальной проверке модели на датасете YAMBDA.
